\documentclass[12pt]{report}
\usepackage[top=2.3cm,bottom=2.5cm,left=2.5cm,right=2.5cm]{geometry}
\usepackage[T1]{fontenc}
\usepackage{mathpazo}
\linespread{1.03}
\usepackage{microtype}
\usepackage{graphicx}
\usepackage{booktabs}
\usepackage{enumitem}
\setlist{itemsep=0.15em,topsep=0.35em,parsep=0pt}
\usepackage{array}
\usepackage{float}
\usepackage{longtable}
\usepackage{tabularx}
\usepackage{amsmath}
\usepackage{amssymb}
\setlength{\parindent}{1.4em}
\setlength{\parskip}{0pt}
\usepackage[hidelinks]{hyperref}

\begin{document}

\begin{titlepage}
\thispagestyle{empty}
\centering
\vspace*{2cm}
{\LARGE Data Structures}\\[1cm]
{\Large CSC 301}\\[1cm]
{\Large Lecture Notes}\\[1.5cm]
Computer Science\\
Osun State University\\
2025/2026\\[2cm]
Prepared by\\
Tunde Sardine
\end{titlepage}

\pagestyle{plain}
\chapter*{Contents}
\addcontentsline{toc}{chapter}{Contents}
\vspace{0.5em}
\noindent
Chapter 1: discuss the appropriate use of built-in data structures; \dotfill\ 3\\[0.6em]
Chapter 2: apply object-oriented concepts (inheritance, polymorphism, design patterns, etc.) in \dotfill\ 12\\[0.6em]
Chapter 3: implement various data structures and their algorithms, and apply them in implementing \dotfill\ 20\\[0.6em]
Chapter 4: choose the appropriate data structure for modelling a given problem; \dotfill\ 30\\[0.6em]
Chapter 5: analyse simple algorithms and determine their efficiency using big-O notation; and \dotfill\ 39\\[0.6em]
Chapter 6: apply the knowledge of data structures to other application domains like data compression \dotfill\ 50\\[0.6em]

\end{document}
